\section{Procedimiento}
\label{sec:procedimiento}

Cada etapa tiene una puerta que puede fallar, y en cada caso se nombra la entrada
sobre la que falla. Una puerta que no puede fallar no es una puerta.

\subsection{Cambios de instrumentación, en orden}

Los once cambios siguen todos la plantilla de tres partes de la
sección~\ref{sec:proposito}: una fila, un guardia que se niega, y dos pruebas.

\begin{center}
\small
\begin{xltabular}{\linewidth}{@{}c L{5.4cm} L{2.2cm} X@{}}
\toprule
\# & cambio & coste & se niega a\\
\midrule
\endfirsthead
\toprule
\# & cambio & coste & se niega a\\
\midrule
\endhead
1  & filtrar por estado la lectura de umbrales & 1 línea & que una declaración retirada gobierne algo\\
2  & publicar el umbral y el resultado de la comparación junto al nodo & pequeño & nada. Hace visible una negativa que hoy se descarta\\
3  & el marco se \emph{designa}, no es la fila más reciente & pequeño & un marco elegido por orden de inserción\\
4  & promover el veredicto emparejado a una tabla tipada & migración y backfill de 9{,}138 lecturas & que borrar un trabajo borre la única evidencia\\
5  & transportar la sexta palabra del vocabulario & pequeño & que un nulo con poder se lea como una pregunta no hecha\\
6  & la declaración como espina de atribución, con clave foránea desde el resultado & migración y 315 filas & un número publicado sin declaración detrás\\
7  & materializar \cod{th\_step} y el estadístico, y añadirlos al sello & migración y re-sellado & comparar dos rejillas distintas, o dos estadísticos distintos\\
8  & retirar o restringir la etiqueta de ventana del sello & re-sellado & que una etiqueta del arnés parta un marco real\\
9  & negarse a comparar arms con bancos de donantes desiguales & un campo y un guardia & el confundido de la sección~\ref{sec:parametros}\\
10 & hacer que los bloqueos estructurales no se lean como evidencia & pequeño & reportar un veredicto que ningún experimento puede mover\\
11 & ligar la ablación de rasgos al esquema versionado & preparado & una ablación que retire familias en lugar de columnas\\
\bottomrule
\caption{Los once cambios de instrumentación. Todos se ejecutan en la máquina que
posee el estado.}
\end{xltabular}
\end{center}

\subsection{Etapas y puertas}

\begin{description}[leftmargin=0pt,itemsep=0.5em,font=\normalfont\bfseries]
  \item[S0. Designar el marco.] Productos: la ventana de selección designada, la
        variante designada y el punto de validación interna. Puerta \textbf{G0}: la
        ventana designada resuelve a una release menor o igual que 227, y la
        comprobación de comparabilidad deja de negarse. \emph{Falla sobre} una
        designación por encima de 227, una ventana cuyo annotation set no existe, o
        una designación que nombre dos evaluation sets a la vez.

  \item[S1. Sacar la evidencia del registro de eventos.] Puerta \textbf{G1}: en una
        transacción que se revierte, intentar borrar un trabajo de comparación debe
        \emph{levantar} por la restricción de la tabla de veredictos, y el recuento
        restaurado debe reproducir las cifras medidas. \textbf{Nada más corre antes
        de G1}, porque toda etapa posterior añade trabajos cuyo borrado se llevaría
        la evidencia.

  \item[S2. Re-sellar bajo los campos corregidos.] Puerta \textbf{G2}: ningún sello
        contiene dos rejillas de umbral ni dos estadísticos. \emph{Falla sobre el
        registro actual}, donde los cuatro sellos contienen las dos rejillas.

  \item[S3. El sustrato, sobre bancos iguales.] Puerta \textbf{G3}: la comparación
        se niega si los bancos difieren, y la negativa es visible en la respuesta.
        \emph{Falla sobre} el censo actual de embeddings.

  \item[S4. El corte de evaluación.] Puerta \textbf{G4}: retirar los resultados con
        corte cambia el número de niveles de algún nodo. \emph{Falla hoy}, porque
        ningún nodo lee el campo del corte.

  \item[S5. Canalización con ajuste cruzado.] Puerta \textbf{G5}: el veredicto
        publicado es fuera de pliegue, y la comparación se niega cuando la partición
        de ajuste y la de evaluación se cortan.

  \item[S6. La medida de validación, una sola vez.] Precondición: todas las puertas
        anteriores superadas y el campeón congelado. Puerta \textbf{G6}: antes de
        calcular nada, el guardia de fuga debe \emph{demostrar una negativa} sobre
        ese marco y después permitir el paso al campeón.
\end{description}

\begin{aviso}{Por qué G6 es la puerta que sostiene todo el protocolo}
Sobre la ventana adoptada, $220 \rightarrow 227$, el guardia de fuga \emph{no puede
levantar}: las tres codificaciones ajustadas que existen declaran corte en la
release 220, y la condición de negativa exige que el corte caiga estrictamente
dentro de la ventana. El guardia está vivo, se invoca, y de forma demostrable no dispara
para ninguna fila de esta base sobre esta ventana. Su única negativa demostrada está
en una prueba sobre otra ventana. La mitad \emph{actúa} de ese guardia queda sin
evidenciar hasta S6, y ahí es donde tiene que demostrarse.
\end{aviso}

\subsection{Reparto entre máquinas}

Una máquina posee el estado: base de datos, cola de trabajos, almacén de objetos, la
interfaz y todo artefacto que la campaña produce. Las demás son nodos de cómputo sin
estado: se suman a la cola, calculan y publican de vuelta. Un nodo de cómputo puede
reiniciarse sin aviso y eso es inocuo por diseño, porque el estado y la cola están en
la otra máquina y una corrida continúa a menor rendimiento hasta que el nodo vuelve.

\begin{center}
\begin{tabularx}{\linewidth}{@{}X l@{}}
\toprule
trabajo & máquina\\
\midrule
codificación (GPU) & nodo de cómputo\\
predicción y recuperación KNN (CPU) & nodo de cómputo\\
evaluación y estratificación & cualquiera\\
comparación emparejada & la que posee el estado\\
ingesta de GOA y ontologías, migraciones, sellado & la que posee el estado, en serie\\
\bottomrule
\end{tabularx}
\end{center}
